\documentclass[11pt]{article}
\usepackage[margin=0.75in]{geometry}
\usepackage{fontspec}
\usepackage{setspace}
\usepackage{hyperref}
\hypersetup{colorlinks=true, urlcolor=blue, linkcolor=black}
\pagestyle{empty}
\setlength{\parskip}{0.65em}
\setlength{\parindent}{0pt}
\setstretch{0.94}
\setmainfont{Times New Roman}

\begin{document}

\begin{flushleft}
Ke Cheng HSU (Clio) \\
Department of Economics, The Hong Kong University of Science and Technology \\
Clear Water Bay, Kowloon, Hong Kong $\cdot$ +852 6840 9753 \\
\href{mailto:khsu@connect.ust.hk}{khsu@connect.ust.hk} $\cdot$ \href{https://cliohsu1997.github.io/}{cliohsu1997.github.io}
\end{flushleft}

18 November 2026

\vspace{0.12em}
Hiring Committee \\
Korea Institute for Industrial Economics \& Trade (KIET) \\
Sejong, Republic of Korea

\vspace{0.12em}
Dear Members of the Hiring Committee,

I am writing to apply for the Associate Research Fellow (Doctoral) position at the Korea Institute for Industrial Economics and Trade. I am a Ph.D.\ candidate in Economics at the Hong Kong University of Science and Technology. I expect to graduate in June 2027 and will have completed the Ph.D.\ by the time employment begins. My fields are Empirical Industrial Organization and Digital Economics. I study firm and consumer decisions under uncertainty---return policy, product design, inventory, and pricing---and quantify how market rules and policy instruments change welfare, which maps directly to KIET's work on industrial organization, trade-facing digital markets, and economic policy.

My job market paper, ``Designing Return Policy in E-Commerce: A Trade-off between Purchase and Return Efficiency,'' studies a core industrial-policy margin for online retail: how a return rule shapes purchases and returns. A stricter policy curbs inefficient returns after delivery, but the lower return rate makes a higher product price more attractive for the seller and can exclude efficient purchases. Using transaction data from a Taobao apparel retailer, I estimate a structural model of purchase, return, and whether to buy return-shipping insurance under monopoly pricing across six categories. Automatic coverage of return shipping performs better when the margin, net scrap value, and uncertainty are high, because it expands efficient purchases. Optional coverage performs better when those attributes are low, because it curbs excessive returns. A menu that pairs a discounted price with a medium restocking fee balances the two margins and performs better than both insurance policies for most categories. The paper yields ranked policy recommendations that a research institute can communicate to industry and government.

A second working paper, ``Product Design with Heterogeneous Search Costs'' (with Pak Hung Au), studies how firms choose product design when consumers differ in search costs. Search costs shape which products consumers search. Consumers with low search costs seek niche products, and consumers with high search costs seek standardized products. A fall in search costs can raise prices when this shift in which product consumers search outweighs the effect that encourages more search. I also have two ongoing projects. With Kohei Kawaguchi and Kei Kawai, I use Japanese bookstore data to study how resale price maintenance encourages retailers to hold inventory in low-demand states, at an environmental cost when unsold copies are returned and discarded. With Ce Bian and Liang Zhao, I study pass-through when products can be returned, and show that a tax can lower the equilibrium price and raise welfare by reducing the return rate.

At HKUST I was a teaching assistant from 2020 to 2025 for undergraduate and Ph.D.\ courses, with an average evaluation of 4.72/5, and in 2026 I designed and lectured the mathematics camp for the MSc in Economics, receiving 4.53/5. I write and present clearly in English and Mandarin and am prepared to contribute to KIET's document screening, interview, and on-site presentation stages with my job market paper.

I submit this letter, my CV, and my job market paper through JOE. Pak Hung Au, Kohei Kawaguchi, and Kei Kawai will submit their letters through JOE. You can reach me at \href{mailto:khsu@connect.ust.hk}{khsu@connect.ust.hk} or +852 6840 9753.

\vspace{0.25em}
Sincerely,

\vspace{0.25em}
Ke Cheng HSU \\
Ph.D.\ Candidate, Department of Economics, HKUST

\end{document}
